\documentclass[11pt]{report}
\usepackage[a4paper,margin=1in]{geometry}
\usepackage{amsmath,amssymb}
\usepackage{booktabs}
\usepackage{longtable}
\usepackage{array}
\usepackage[colorlinks=true,linkcolor=blue,urlcolor=blue]{hyperref}
\usepackage{fancyhdr}
\usepackage{xcolor}
\usepackage{titlesec}
\usepackage{enumitem}

\definecolor{hbblue}{HTML}{2A4D8F}
\definecolor{warnred}{HTML}{9C2B2B}
\definecolor{boxgrey}{HTML}{F4F4F6}

\titleformat{\chapter}[display]{\bfseries\Large\color{hbblue}}{\chaptertitlename\ \thechapter}{10pt}{\huge}
\pagestyle{fancy}
\fancyhf{}
\fancyhead[R]{\small\itshape Single-Pinhole logMAR Acuity Protocol}
\fancyfoot[C]{\thepage}

\newsavebox{\kbsave}
\newsavebox{\wbsave}
\newenvironment{keybox}[1]%
{\par\vspace{8pt}\begin{lrbox}{\kbsave}\begin{minipage}{0.88\textwidth}%
{\bfseries\color{hbblue}#1}\par\vspace{4pt}\ignorespaces}%
{\end{minipage}\end{lrbox}%
\begin{center}\setlength{\fboxrule}{1.2pt}\setlength{\fboxsep}{8pt}%
\fcolorbox{hbblue}{boxgrey}{\usebox{\kbsave}}\end{center}\vspace{8pt}\par}
\newenvironment{warnbox}[1]%
{\par\vspace{8pt}\begin{lrbox}{\wbsave}\begin{minipage}{0.88\textwidth}%
{\bfseries\color{warnred}#1}\par\vspace{4pt}\ignorespaces}%
{\end{minipage}\end{lrbox}%
\begin{center}\setlength{\fboxrule}{1.2pt}\setlength{\fboxsep}{8pt}%
\fcolorbox{warnred}{white}{\usebox{\wbsave}}\end{center}\vspace{8pt}\par}

\title{\textbf{\Huge Single-Pinhole Acuity}\\[0.3cm]
\textcolor{hbblue}{\huge logMAR Measurement Protocol}\\[0.5cm]
\large\itshape Model predictions read from \texttt{acuity\_sweep\_stage4.csv},\\
and the experiment that tests them}
\author{}
\date{IYPT 2027 Problem 1}

\begin{document}
\maketitle
\tableofcontents

\chapter*{Provenance of every number in this document}
\addcontentsline{toc}{chapter}{Provenance}

\begin{warnbox}{Read this first}
Every predicted acuity value in this document is read directly from
\texttt{new\_sweeps/acuity\_sweep\_stage4.csv}, column \texttt{logmar}, produced by the
\texttt{mtf10\_vsotf\_brightness\_adjusted} model. No value here is estimated, interpolated, or
assumed. Where a quantity is \emph{not} available from the model --- most importantly the
bare-eye baseline --- this document says so explicitly and does not supply a substitute.
\end{warnbox}

\begin{warnbox}{This model failed its one external validation check --- read before quoting any number below}
The full acuity model (every number in this document) was tested against Wang et al.\ (2009)'s real
clinical pinhole-acuity data, in \texttt{stage4\_human\_vision/RESULTS.md}: $R^2 = -6.39$ to $-2.84$
across $M=0,2,5$, $p\approx0.50$. Negative $R^2$ means worse than guessing the mean; $p\approx0.50$
means no significant relationship was detected. Every one of the 15 tested points showed the same
bias direction --- the model over-predicts blur at every point.

The scene-luminance assumption was tested and ruled out as the cause: correcting to Wang 2009's own
reported illumination made $R^2$ \emph{worse}, not better. The project's own leading diagnosis is an
uncalibrated formula combined with a small ($n=5$), non-monotonic clinical dataset --- not
necessarily wrong optics. Full reasoning in \texttt{stage4\_human\_vision/300TD\_KNEE\_NOTES.txt}.

\textbf{This is exactly why the experiment below matters.} Every number in this document is a
self-consistent, traceable model output --- not a validated clinical prediction. The measurement you
are about to run is calibration data this model currently has none of.
\end{warnbox}

\chapter{What the model predicts, and what it does not}

\section{What it predicts}

The model predicts \textbf{acuity through a pinhole of stated diameter, at stated myopia}. That is
one endpoint of the experiment.

\section{What it does not predict}

The model does \textbf{not} predict uncorrected bare-eye acuity. This is not an oversight --- it is a
direct consequence of the declared scope. Uncorrected acuity in a real eye is set by defocus
\emph{plus} higher-order aberrations, accommodation, and neural factors, all of which are explicitly
excluded from this project's model. The simulation contains a pinhole in every configuration it can
evaluate; there is no aperture-free case in the sweep.

\begin{keybox}{The consequence for your experimental claim}
You cannot state a predicted \emph{improvement}. You can state a predicted \emph{endpoint}, and
measure the improvement against a \emph{measured} baseline. The defensible sentence is:

\emph{``The model predicts logMAR $+0.33$ through a 400\,\textmu m aperture at $M=8$\,D. We measured
logMAR $X$ through that aperture, against a measured bare-eye baseline of $Y$.''}

Anyone claiming a predicted $\Delta$ from this model is asserting a bare-eye number the model never
produced.
\end{keybox}

\chapter{Predicted acuity for the fabricated designs}

Each row pairs the gaze-weighted winning aperture for that myopia level with the model's predicted
acuity at exactly that aperture and myopia.

\begin{table}[h!]
\centering
\begin{tabular}{cccccc}
\toprule
$M$ (D) & Aperture (\textmu m) & MAR (arcmin) & logMAR & Snellen & Retinal illum.\ (Td) \\
\midrule
0 & 2000 & 0.45 & $-0.344$ & 20/9.1 & 31416 \\
1 & 1250 & 0.76 & $-0.120$ & 20/15.2 & 12272 \\
2 & 750 & 1.19 & $+0.076$ & 20/23.8 & 4418 \\
3 & 750 & 1.29 & $+0.112$ & 20/25.9 & 4418 \\
4 & 600 & 1.52 & $+0.181$ & 20/30.3 & 2827 \\
5 & 600 & 1.74 & $+0.240$ & 20/34.8 & 2827 \\
6 & 500 & 1.81 & $+0.259$ & 20/36.3 & 1963 \\
8 & 400 & 2.14 & $+0.331$ & 20/42.8 & 1257 \\
\bottomrule
\end{tabular}
\caption{Predicted acuity through each fabricated design. Apertures from
\texttt{gaze\_weighted\_winner\_per\_M.csv}; acuity from \texttt{acuity\_sweep\_stage4.csv}.}
\end{table}

\begin{keybox}{An independent consistency check, worth stating at defence}
These winning apertures were selected by the \emph{gaze-weighted geometric} evaluator --- viability,
ghost margin, blind fraction --- with no acuity term in its objective. Yet at every myopia level from
0 to 8\,D, the selected aperture is \textbf{also} the acuity-optimal aperture in the independent
Stage-4 sweep. Two evaluators built on different objectives agree on all eight designs. That
agreement was not enforced and is not a coincidence worth passing over.
\end{keybox}

\chapter{The strongest testable prediction you have}
\label{ch:flatness}

\section{Acuity through a fixed small aperture is nearly independent of myopia}

Holding the aperture at 400\,\textmu m and varying only myopia, the model gives:

\begin{table}[h!]
\centering
\begin{tabular}{ccc}
\toprule
$M$ (D) & logMAR & Snellen \\
\midrule
0 & $+0.341$ & 20/43.9 \\
1 & $+0.337$ & 20/43.4 \\
2 & $+0.333$ & 20/43.0 \\
3 & $+0.329$ & 20/42.7 \\
4 & $+0.327$ & 20/42.4 \\
5 & $+0.325$ & 20/42.3 \\
6 & $+0.325$ & 20/42.2 \\
8 & $+0.331$ & 20/42.8 \\
\bottomrule
\end{tabular}
\caption{Predicted acuity through a fixed 400\,\textmu m aperture across the full myopia range.
Total spread: 0.016 logMAR.}
\end{table}

Across eight dioptres of myopia the predicted acuity varies by \textbf{0.016 logMAR} --- roughly a
twentieth of one chart line, and far below the $\pm0.1$ logMAR repeatability of human acuity
testing. At 400\,\textmu m the aperture, not the refractive error, sets the acuity.

\begin{keybox}{Why this is the experiment to run}
This prediction is \textbf{falsifiable, sharp, and requires no knowledge of anyone's exact
prescription}. Recruit subjects across a range of myopia. If the model is right, all of them read
approximately the same line through the same 400\,\textmu m pinhole --- while their bare-eye acuity
differs enormously.

That is a far stronger experimental design than measuring one subject and comparing to one number,
because the \emph{shape} of the result carries the claim. A flat line across subjects of very
different refractive error is difficult to produce by accident.
\end{keybox}

\section{Why the flatness happens}

Two effects nearly cancel. As myopia increases, the defocus blur circle grows --- but the pinhole is
already narrow enough that the accepted cone is set by the aperture, not by the eye's focusing error.
The residual defocus contribution is small compared with the diffraction floor the 400\,\textmu m
aperture imposes. The predicted acuity therefore tracks the diffraction limit of the aperture, which
does not depend on $M$ at all.

The flatness breaks at larger apertures, where defocus reasserts itself. Compare 2000\,\textmu m:
logMAR $-0.344$ at $M=0$ but $+1.330$ at $M=8$ --- a spread of 1.67 logMAR, a hundred times larger.
That contrast is itself a second testable prediction, and a good slide.

\chapter{Building the chart}

\section{logMAR, defined}

logMAR is $\log_{10}$ of the minimum resolvable angle in arcminutes. A one-arcminute detail gives
logMAR $=0$, conventionally 20/20. Positive is worse, negative is better. One chart line is 0.1
logMAR by construction, so a line's letters are $10^{0.1} = 1.259\times$ the size of the next line down.

\begin{equation}
\text{logMAR} = \log_{10}(\text{MAR in arcmin}), \qquad
\text{Snellen denominator} = 20 \times 10^{\text{logMAR}}
\end{equation}

\section{Letter sizing}

A Snellen letter is five times the stroke width, and the stroke subtends the MAR. So at viewing
distance $D$, the letter height is
\begin{equation}
h = 5 \times \text{MAR}_{\text{arcmin}} \times \frac{\pi}{180\times 60} \times D
= 1.4544\times10^{-3}\,\text{MAR}_{\text{arcmin}} \times D
\end{equation}

At $D = 4.0$\,m:

\begin{table}[h!]
\centering
\begin{tabular}{ccccc}
\toprule
logMAR & MAR (arcmin) & Snellen & Letter height (mm) & Stroke (mm) \\
\midrule
$+0.7$ & 5.01 & 20/100 & 29.2 & 5.83 \\
$+0.6$ & 3.98 & 20/79.6 & 23.2 & 4.63 \\
$+0.5$ & 3.16 & 20/63.2 & 18.4 & 3.68 \\
$+0.4$ & 2.51 & 20/50.2 & 14.6 & 2.92 \\
$+0.3$ & 2.00 & 20/39.9 & 11.6 & 2.32 \\
$+0.2$ & 1.58 & 20/31.7 & 9.2 & 1.84 \\
$+0.1$ & 1.26 & 20/25.2 & 7.3 & 1.46 \\
$0.0$ & 1.00 & 20/20 & 5.8 & 1.16 \\
$-0.1$ & 0.79 & 20/15.9 & 4.6 & 0.92 \\
$-0.2$ & 0.63 & 20/12.6 & 3.7 & 0.73 \\
\bottomrule
\end{tabular}
\caption{Letter dimensions at 4.0\,m viewing distance. Scale linearly for other distances.}
\end{table}

\begin{keybox}{Why 4 metres, and why this range}
The predictions cluster between logMAR $-0.34$ and $+0.34$, so the chart must span at least
$-0.2$ to $+0.7$ to bracket both the predicted pinhole endpoints and plausible bare-eye baselines
for mild myopes. At 4\,m the smallest stroke is 0.92\,mm --- printable on a domestic laser printer.
At 2\,m it would be 0.46\,mm, at the edge of print fidelity, and printer dot structure would start
to limit the measurement rather than the eye.
\end{keybox}

\section{Construction}

Five letters per line, Sloan set (C, D, H, K, N, O, R, S, V, Z --- equal legibility by design), even
spacing of one letter width. Print black on matte white; glossy paper produces specular glare that
confounds the low-transmission conditions. Verify the print scale by measuring a printed reference
bar with callipers before use, exactly as with the etch masks.

\chapter{The confounder that will ruin this measurement}

\section{Retinal illuminance is not a footnote}

The predicted-acuity table carries a Troland column for a reason. At $M=8$ the 400\,\textmu m design
delivers 1257\,Td against 31416\,Td for the 2000\,\textmu m case --- a factor of 25. The model's
\texttt{mtf10\_vsotf\_brightness\_adjusted} predictions already include a brightness adjustment, but
they assume a specific scene luminance (10\,000\,cd/m$^2$, a bright outdoor surface).

\begin{warnbox}{If you test indoors, the model's assumption is violated}
A typical indoor chart under room lighting is 100--300\,cd/m$^2$, some 30--100$\times$ dimmer than
assumed. Through a 400\,\textmu m aperture at 1\% transmission, the retinal illuminance can fall low
enough that acuity becomes photon-limited rather than optics-limited --- and you will measure a
brightness effect while believing you measured a diffraction effect.

\textbf{Mitigation: illuminate the chart hard.} Direct lamps onto the chart to raise its luminance as
high as you can and \emph{measure it} with a light meter (a phone app with a calibrated lux reading is
adequate to within a factor of two, which is enough to know whether you are in the right regime).
Record the value. If you cannot reach high luminance, state that the bare-eye and pinhole conditions
were measured at different effective retinal illuminance, and that the comparison is therefore not
purely optical.
\end{warnbox}

\section{Accommodation}

At 4\,m a myopic eye cannot accommodate to bring the chart into focus --- accommodation only adds
power, moving focus nearer, which is the wrong direction. That is convenient: it means the defocus
you are testing is stable. \textbf{Do not test at short distances}, where accommodative effort varies
between trials and between conditions.

\section{Pupil diameter}

The model assumes a 4\,mm pupil. Under the bright illumination this protocol requires, real pupils
constrict to 2--3\,mm. A smaller pupil improves bare-eye acuity by itself --- a natural pinhole
effect --- which \emph{reduces} the apparent benefit of the artificial aperture. Photograph the eye
against a mm scale under test lighting, or at minimum record the lighting conditions so the
direction of this bias is documented.

\section{Learning effect}

Subjects memorise charts. Use at least two different letter arrangements and alternate them, or
randomise the letter order between repeats. Otherwise later conditions score artificially well.

\section{The ranked list of things that will go wrong}

\begin{enumerate}[leftmargin=1.2cm]
\item Chart too dim, so you measure photon noise instead of optics.
\item Chart printed at the wrong scale --- verify with callipers.
\item Pupil constricted under bright light, shrinking the measured benefit.
\item Subject memorising the chart.
\item Pinhole not centred on the pupil, so the subject sees through its edge or misses it entirely.
\item $d_{\text{gap}}$ varying between conditions because the mount is hand-held.
\end{enumerate}

\chapter{Protocol}

\section{Fixed conditions --- record once per session}

\begin{longtable}{p{5.5cm}p{4.0cm}p{4.5cm}}
\toprule
\textbf{Quantity} & \textbf{Value} & \textbf{How obtained} \\
\midrule
\endhead
Viewing distance & \rule{3cm}{0.4pt}\,m & tape measure, $\pm0.02$\,m \\
Chart luminance & \rule{3cm}{0.4pt}\,cd/m$^2$ or lux & light meter / phone app \\
Chart version(s) used & \rule{3cm}{0.4pt} & A / B alternation \\
$d_{\text{gap}}$ (mount to cornea) & \rule{3cm}{0.4pt}\,mm & spacer, fixed \\
Room lighting & \rule{3cm}{0.4pt} & description \\
Date / time & \rule{3cm}{0.4pt} & \\
\bottomrule
\end{longtable}

\section{Per-subject record}

\begin{longtable}{p{5.5cm}p{4.0cm}p{4.5cm}}
\toprule
\textbf{Quantity} & \textbf{Value} & \textbf{Source} \\
\midrule
\endhead
Subject ID & \rule{3cm}{0.4pt} & anonymised \\
Refractive error, tested eye & \rule{3cm}{0.4pt}\,D & optometrist Rx if available \\
Rx date & \rule{3cm}{0.4pt} & \\
Astigmatism (cyl) & \rule{3cm}{0.4pt}\,D & note if $>1$\,D \\
Eye tested & L / R & occlude the other \\
Pupil diameter under test light & \rule{3cm}{0.4pt}\,mm & photo vs.\ scale, or note \\
\bottomrule
\end{longtable}

\begin{keybox}{If prescriptions are unavailable}
Chapter~\ref{ch:flatness}'s flatness prediction does not need them. Rank subjects qualitatively
(``no glasses'', ``mild'', ``strong'') and test whether pinhole acuity is flat across that ranking
while bare-eye acuity is not. Report refractive error as unmeasured. A qualitative ordering with an
honest caveat is worth more than a fabricated dioptre value.
\end{keybox}

\section{Procedure}

\begin{enumerate}[leftmargin=1.2cm]
\item Occlude the untested eye. Keep it open under the occluder.
\item Seat the subject at the fixed distance, head steady.
\item Present the chart. Subject reads from the largest line downward, aloud, letter by letter.
\item Score a line as read if at least 3 of 5 letters are correct. The threshold is the smallest such line.
\item Record the line's logMAR.
\item Rest 30\,s. Change chart version.
\item Repeat for the next condition in randomised order.
\end{enumerate}

\textbf{Conditions:} (1) bare eye, (2) single 400\,\textmu m pinhole, (3) fabricated array.
Three repeats each, order randomised across all nine trials.

\section{Recording table}

\begin{longtable}{cccccc}
\toprule
Trial & Subject & Condition & Chart & Threshold line & logMAR \\
\midrule
\endhead
1 & & & & & \\
2 & & & & & \\
3 & & & & & \\
4 & & & & & \\
5 & & & & & \\
6 & & & & & \\
7 & & & & & \\
8 & & & & & \\
9 & & & & & \\
\bottomrule
\end{longtable}

\chapter{Analysis}

\section{Per condition}

Mean and standard deviation of the three repeats. Human acuity repeatability is typically
$\pm0.1$ logMAR, so a standard deviation much above 0.15 indicates a procedural problem --- most
often inconsistent pinhole centring --- rather than a real effect.

\section{Test 1 --- the endpoint comparison}

For each subject, compare measured pinhole acuity against the model's prediction at that aperture
and that subject's myopia (Chapter~3 table).

\begin{equation}
\text{residual} = \text{logMAR}_{\text{measured}} - \text{logMAR}_{\text{predicted}}
\end{equation}

Report the residual, signed. A positive residual means the eye did worse than predicted. Interpretation:

\begin{itemize}[leftmargin=1.2cm]
\item $|{\text{residual}}| < 0.15$: agreement within measurement repeatability.
\item Residual $> +0.3$ consistently: a confounder is active. Check chart luminance first --- the model
assumes bright outdoor conditions.
\item Residual $< -0.3$: the eye outperforms the model. The model omits neural processing and
excludes higher-order aberrations, both of which can favour the real eye.
\end{itemize}

\section{Test 2 --- the flatness test (the strong one)}

Plot measured pinhole logMAR against myopia (or qualitative rank) for all subjects. The model
predicts a spread of 0.016 logMAR across 0--8\,D --- effectively a horizontal line. On the same axes
plot bare-eye logMAR, which should fall steeply with increasing myopia.

\begin{keybox}{What each outcome licenses you to say}
\textbf{Pinhole flat, bare-eye steep:} the central claim of the whole project, measured. \emph{``Through
a 400\,\textmu m aperture, acuity is set by the aperture rather than by refractive error --- measured
across subjects spanning X dioptres.''}

\textbf{Pinhole also falls with myopia:} the model is wrong in a specific, interesting way, and the
slope tells you how wrong. This is a publishable-grade negative result, not a failure.

\textbf{Both flat:} your subjects are not myopic enough to separate the conditions, or the chart is
too easy. Extend the chart to smaller lines and recruit a stronger myope.
\end{keybox}

\section{Test 3 --- array versus single hole}

The sparse array is designed so that only one aperture illuminates the pupil at a time. Two
consequences follow, both testable:

\begin{enumerate}[leftmargin=1.2cm]
\item \textbf{On axis, array acuity should equal single-hole acuity} to within repeatability, since
the eye is looking through one hole in both cases.
\item \textbf{Off axis, the array should keep working and the single hole should not.} Rotate the
subject's gaze $10$--$20^\circ$ by moving the chart laterally while the mount stays fixed. The single
hole should go dark; the array should hand over to a neighbouring aperture.
\end{enumerate}

Also record: \textbf{does any subject report a doubled image at any gaze angle?} The model predicts
none, at any myopia, because no ghost pair ever forms. An observed ghost would falsify that
prediction and is worth more than confirming it.

\chapter{Reporting}

\section{The claim, written honestly}

\begin{quote}
``Through a 400\,\textmu m aperture at $d_{\text{gap}} = $\,\rule{1cm}{0.4pt}\,mm and
\rule{1cm}{0.4pt}\,m viewing distance, $N$ subjects spanning \rule{1cm}{0.4pt}--\rule{1cm}{0.4pt}\,D
of myopia read logMAR \rule{1cm}{0.4pt} $\pm$ \rule{1cm}{0.4pt}, against bare-eye acuity ranging from
\rule{1cm}{0.4pt} to \rule{1cm}{0.4pt}. Our model predicts $+0.33 \pm 0.01$ across that same range;
this model has not previously been validated against real clinical data, and our one prior validation
attempt (against Wang et al.\ 2009) found no significant relationship ($R^2<0$). This measurement is
that calibration.''
\end{quote}

\section{What must accompany it}

\begin{itemize}[leftmargin=1.2cm]
\item Chart luminance, measured, with a statement of whether it met the model's assumed regime.
\item Pupil diameter, measured or noted as unmeasured.
\item Sample size and how refractive error was established.
\item Every subject's raw data --- do not report only the mean.
\item Any subject excluded, and why, decided before analysis.
\end{itemize}

\begin{warnbox}{Do not report a predicted improvement}
The model produces no bare-eye baseline. Any $\Delta$ you quote must have a \emph{measured} baseline
on both sides. Quoting a predicted improvement means quoting a bare-eye number the model never
generated --- which is exactly the kind of unsourced figure this project has otherwise been careful
to avoid.
\end{warnbox}

\end{document}
